\section*{Chapter \ref{ch:iv:mental-arithmetic} --- Mental Arithmetic}

\begin{solution}{iq:iv:mental-arithmetic:1}
$47 \times 53 = 50^2 - 3^2 = 2\,500 - 9 = 2\,491$.
\iqlookfor{spotting the difference of squares at once.}
\end{solution}

\begin{solution}{iq:iv:mental-arithmetic:2}
$96 \times 104 = 100^2 - 4^2 = 9\,984$. $98 \times 97 = 100(100 - 2 - 3) + (-2)(-3) = 9\,500 + 6 = 9\,506$.
\iqlookfor{the near-hundred rewrite, including negative offsets.}
\end{solution}

\begin{solution}{iq:iv:mental-arithmetic:3}
$65^2 = 100 \times 6 \times 7 + 25 = 4\,225$. $75^2 - 65^2 = (75 - 65)(75 + 65) = 10 \times 140 = 1\,400$.
\iqlookfor{the squares-ending-in-5 rule and factoring a difference of squares rather than computing both.}
\end{solution}

\begin{solution}{iq:iv:mental-arithmetic:4}
12.5\% is $\tfrac18$: $360/8 = 45$. 36\% of 25 is 25\% of 36: 9.
\iqlookfor{percentages as fractions, and the symmetry of a percentage.}
\end{solution}

\begin{solution}{iq:iv:mental-arithmetic:5}
$\tfrac37 = 3 \times 0.142857 \approx 0.429$. $\tfrac5{13} = 5 \times 0.0769 \approx 0.385$.
\iqlookfor{a reciprocal table used as multiplication.}
\end{solution}

\begin{solution}{iq:iv:mental-arithmetic:6}
$\sqrt{49 + 1} \approx 7 + \tfrac1{14} \approx 7.0714$, with error below $1/(8 \times 343) \approx 0.00036$; the
true value is $7.0711$, so the estimate is $0.0003$ high, inside the bound. (A second-order term, $-1/(8 \times
343)$, corrects it to $7.0711$.)
\iqlookfor{linearisation around a known square and a stated error bound.}
\end{solution}

\begin{solution}{iq:iv:mental-arithmetic:7}
The rule gives 72, 12 and 3 years; the exact values are 69.7, 11.9 and 3.22. The 24\% answer is furthest in
relative terms (about 7\% short); the 1\% answer is about 3.4\% long. The rule overestimates below about 7.85\%
and underestimates above it (\cref{prop:iv:mental-arithmetic:rule72}).
\iqlookfor{the rule, the exact formula, and knowing the direction of the error.}
\end{solution}

\begin{solution}{iq:iv:mental-arithmetic:8}
3.2 basis points of 250 million is $3.2 \times 250 \times 100 = 80\,000$ a day; over 252 days, $80\,000 \times
252 = 20.16$ million.
\iqlookfor{basis points of millions as hundreds, and a clean multiplication by 252 ($250 + 2$).}
\end{solution}

\begin{solution}{iq:iv:mental-arithmetic:9}
Residues: $3\,847 \to 22 \to 4$, $29 \to 2$, product residue 8; $111\,563 \to 17 \to 8$: it passes, and it is
right. $111\,653$ has the same digits, so the same residue: it passes too, and it is wrong. \omterm{def:iv:mental-arithmetic:nines}{Casting out nines}
cannot see a transposition or any error by a multiple of 9; add the last-digit check ($7 \times 9$ ends in 3,
which both pass) and a magnitude check, or recompute one partial product.
\iqlookfor{the check and, more important, its blind spot.}
\end{solution}

\begin{solution}{iq:iv:mental-arithmetic:10}
$1.2 \times 0.8 = 0.96$: 4\% below the start. After a 25\% rise the price is $\tfrac54$ of the start, so a fall of
$\tfrac15$, 20\%, brings it back: $1.25 \times 0.8 = 1$.
\iqlookfor{multiplicative returns and the asymmetry of percentage moves.}
\end{solution}

\begin{solution}{iq:iv:mental-arithmetic:11}
$\ln 1.07 = 0.07 - 0.00245 + 0.000114 - \dots \approx 0.0677$ (to four decimals; the exact value is
$0.06766$). The doubling time is $0.6931/0.06766 \approx 10.24$ years, against 10.29 from the \omterm{def:iv:mental-arithmetic:rule72}{rule of 72}.
\iqlookfor{three terms of the series and a clean division.}
\end{solution}

\begin{solution}{iq:iv:mental-arithmetic:12}
$10 \ln 1.03 = 10(0.03 - 0.00045 + 0.000009) \approx 0.2956$; $e^{0.2956} = e^{0.3} e^{-0.0044} \approx 1.3499
\times 0.9956 \approx 1.344$ (exactly $1.3439$).
\iqlookfor{compounding through logarithms, with $e^{0.3}$ known or derived.}
\end{solution}

\begin{solution}{iq:iv:mental-arithmetic:13}
$0.0625 = \tfrac1{16}$ and $4.8 \times 125 = 600$ (125 is $\tfrac{1000}8$, and $4.8/8 = 0.6$), so the product is
$600/16 = 37.5$.
\iqlookfor{recognising friendly factors before multiplying.}
\end{solution}

\begin{solution}{iq:iv:mental-arithmetic:14}
For annual compounding the exact doubling time is $\ln 2/\ln(1 + r)$ and $\ln(1+r) < r$, so $69.3/r$ is always
too short; a larger numerator compensates, and 72 also has many divisors (2, 3, 4, 6, 8, 9, 12). The rule is
exact where $72/(100r) = \ln 2/\ln(1+r)$, about $r = 7.85\%$: near the rates at which it is usually applied.
For continuous compounding, $\ln 2/r$ is exact, so the rule of 69.3.
\iqlookfor{the source of the rule's error and why its constant is a compromise.}
\end{solution}
